% =====================================================================
%  Course syllabus                                       (Persian edition)
% =====================================================================
\documentclass[11pt,a4paper]{article}
\providecommand\ACNroot{..}
\input{\ACNroot/templates/lecture-fa}

\ACNset{number}{0}
\ACNset{title}{طرح درس}
\ACNset{shorttitle}{طرح درس}

\begin{document}
\begin{center}
  {\ACNheadfont\bfseries\color{acnnavy}\LARGE شبکه‌های ارتباطی پیشرفته\par}
  \vspace{2pt}
  {\ACNheadfont\color{acnteal}\Large طرح درس\par}
  \vspace{2pt}
  {\lr{Advanced Communication Networks -- Course Syllabus}\par}
\end{center}

\section{شرح درس}
این درس شبکه‌های ارتباطی موبایل را از نسل دوم تا نسل ششم دنبال می‌کند. سازوکارهای دسترسی رادیویی، شبکهٔ هسته، سیگنالینگ، جابه‌جایی، امنیت و کیفیت
خدمات در \lr{GSM/GPRS}، \acro{UMTS}، \lr{LTE/EPC} و \lr{5G NR/5GC} و فناوری‌های نرم‌افزارمحوری (\acro{SDN}، \acro{NFV}، معماری خدمت‌محور) که آن‌ها را دگرگون
کردند توضیح داده می‌شوند. هر درس معماری، رویه‌ها و مدل‌های کمّی را با مثال‌های حل‌شده و تمرین‌های درجه‌بندی‌شده ترکیب می‌کند. همهٔ مطالب به انگلیسی و
فارسی با ساختار، نمادگذاری، شکل‌ها و شماره‌گذاری یکسان در دسترس‌اند.

\section{پیامدهای یادگیری}
دانشجویان در پایان درس می‌توانند:
\begin{objectives}
  \item معماری و کارکردهای شبکه‌های نسل دوم تا پنجم را \textbf{شرح دهند} و کارکردها را میان نسل‌ها نگاشت کنند.
  \item رویه‌های اصلی (اتصال/ثبت، احراز هویت، برقراری نشست، دست‌به‌دست، فراخوانی) را در نمودارهای توالی \textbf{دنبال کنند}.
  \item مدل‌های کمّی -- ترافیک و صف، بودجهٔ پیوند، زمان‌بندی، ظرفیت، تأخیر و بار سیگنالینگ -- را برای ابعادگذاری شبکه \textbf{به کار ببرند}.
  \item \acro{SDN}، \acro{NFV} و معماری خدمت‌محور و نقش آن‌ها در برش‌بندی و \acro{QoS} نسل پنجم را \textbf{توضیح دهند}.
  \item مشخصات \lr{3GPP} و \acro{ITU} را \textbf{بخوانند} و فناوری‌های نوظهور نسل ششم را نقادانه \textbf{ارزیابی کنند}.
\end{objectives}

\section{پیش‌نیازها}
احتمال و نظریهٔ صف مقدماتی، سیگنال‌ها و سیستم‌ها، مخابرات دیجیتال (مدولاسیون، ظرفیت شانون) و شبکه‌های رایانه‌ای (\lr{TCP/IP}).

\section{ساختار درس}
۲۶ درس در پنج بخش گروه‌بندی شده‌اند. هر درس سندی مستقل با چهارده بخش است: اهداف، پیش‌نیازها، مقدمه، مفاهیم بنیادی، جزئیات فنی، مدل‌های ریاضی،
معماری، رویه‌ها، مثال‌های حل‌شده، مقایسهٔ فناوری‌ها، نکات کلیدی، تمرین‌ها در سه سطح، پاسخ‌ها و مطالعهٔ بیشتر.

\begin{table}[htbp]
\centering\small
\caption{سرفصل‌های درس.}\label{tab:modules}
\rowcolors{2}{acnmist}{white}
\begin{tabularx}{\linewidth}{@{}r Y l@{}}
\toprule
\rowcolor{white} شماره & درس & نسل \\
\midrule
۱ & جی‌اس‌ام، جی‌پی‌آر‌اس و اج & \lr{2G} \\
۲ & سوئیچ مداری در نسل دوم & \lr{2G} \\
۳ & شبکه‌های هستهٔ نسل دوم و سوم & \lr{3G} \\
۴ & تحرک و تحویل & \lr{4G} \\
۵ & معماری \lr{LTE} نسل چهارم & \lr{4G} \\
۶ & \lr{E-UTRAN}: شبکهٔ دسترسی رادیویی \lr{LTE} & \lr{4G} \\
۷ & زمان‌بندی در \lr{LTE} & \lr{4G} \\
۸ & کانال دسترسی تصادفی (\lr{RACH}) & \lr{4G} \\
۹ & هستهٔ بستهٔ تکامل‌یافته (\lr{EPC}) & \lr{4G} \\
۱۰ & پروتکل \lr{Diameter} & \lr{4G} \\
۱۱ & کارساز مشترکین خانگی (\lr{HSS}) & \lr{4G} \\
۱۲ & موجودیت مدیریت تحرک (\lr{MME}) & \lr{4G} \\
۱۳ & دروازه‌های خدمت‌دهنده و \lr{PDN} (\lr{S-GW} و \lr{P-GW}) & \lr{4G} \\
۱۴ & مقدمه‌ای بر شبکه‌های نسل پنجم & \lr{5G} \\
۱۵ & شبکه‌سازی نرم‌افزارمحور (\lr{SDN}) & \lr{5G} \\
۱۶ & مجازی‌سازی کارکردهای شبکه (\lr{NFV}) & \lr{5G} \\
۱۷ & استانداردسازی \lr{3GPP} & \lr{5G} \\
۱۸ & معماری سرویس‌محور نسل پنجم (\lr{SBA}) & \lr{5G} \\
۱۹ & معماری شبکهٔ هستهٔ نسل پنجم & \lr{5G} \\
۲۰ & رویه‌های سامانهٔ نسل پنجم & \lr{5G} \\
۲۱ & برش‌بندی شبکه & \lr{5G} \\
۲۲ & کیفیت خدمات (\lr{QoS}) در نسل پنجم & \lr{5G} \\
۲۳ & صفحهٔ کاربر نسل پنجم & \lr{5G} \\
۲۴ & صفحهٔ کنترل نسل پنجم & \lr{5G} \\
۲۵ & تکامل شبکه: یک جمع‌بندی بزرگ & \lr{5G} \\
۲۶ & فناوری‌های آیندهٔ نسل ششم & \lr{6G} \\
\bottomrule
\end{tabularx}
\end{table}

\begin{itemize}
  \item \textbf{بخش اول -- مبانی نسل دوم و سوم} (درس‌های ۱ تا ۴).
  \item \textbf{بخش دوم -- دسترسی و هستهٔ نسل چهارم} (درس‌های ۵ تا ۱۳).
  \item \textbf{بخش سوم -- معماری‌های نرم‌افزارمحور} (درس‌های ۱۵ تا ۱۸).
  \item \textbf{بخش چهارم -- سامانه و رویه‌های نسل پنجم} (درس‌های ۱۴ و ۱۹ تا ۲۲).
  \item \textbf{بخش پنجم -- صفحه‌ها و چشم‌انداز} (درس‌های ۲۳ تا ۲۶).
\end{itemize}

\section{برنامهٔ زمانی پیشنهادی}
برای یک نیم‌سال ۱۳ هفته‌ای، هر هفته دو درس به ترتیب شماره؛ هر درس حدود ۶ تا ۸ ساعت مطالعه نیاز دارد (۳ ساعت خواندن و ۳ تا ۵ ساعت تمرین). هفته‌های ۱ و ۲
بخش اول، هفته‌های ۳ تا ۷ بخش دوم و درس ۱۴، هفته‌های ۸ و ۹ بخش سوم، هفته‌های ۱۰ و ۱۱ باقی بخش چهارم و هفته‌های ۱۲ و ۱۳ بخش پنجم را پوشش می‌دهند.

\section{ارزشیابی پیشنهادی}
وزن‌دهی زیر پیشنهادی است و مدرسان می‌توانند آن را تطبیق دهند: مجموعه‌تمرین‌های هفتگی ۳۰٪، آزمون میان‌ترم از بخش‌های اول و دوم ۲۵٪، یک پروژه (مثلاً
ابعادگذاری یا شبیه‌سازی بخشی از یک شبکهٔ نسل پنجم) ۱۵٪ و آزمون پایانی از کل درس ۳۰٪.

\section{منابع}
جزوه‌ها، شکل‌ها و کتاب‌نامهٔ این مخزن؛ اسناد \lr{3GPP} و \acro{ITU} که در هر درس ارجاع داده شده‌اند؛ و کتاب‌های درسی فهرست‌شده در بخش مطالعهٔ بیشتر،
به‌ویژه دالمن و همکاران (\acro{NR})، سزیا و همکاران (\acro{LTE}) و رومر و همکاران (هستهٔ نسل پنجم).

\section{زبان}
دو نسخه هم‌ارزند و دانشجویان می‌توانند به هر یک از دو زبان مطالعه کنند. اختصارهای فنی در هر دو نسخه به شکل لاتین استاندارد خود حفظ شده‌اند.
\end{document}
